\documentclass[11pt, letterpaper]{article}
\usepackage[T1]{fontenc}
\usepackage{amsmath, amssymb, amsthm}
\usepackage[margin=1.12in]{geometry}
\usepackage{lmodern}
\usepackage[hidelinks]{hyperref}
\hypersetup{
    pdftitle={A hereditarily atomic domain without ACCP},
    pdfauthor={AI-authored manuscript}
}
\newtheorem{maintheorem}{Theorem}
\renewcommand{\themaintheorem}{\Alph{maintheorem}}
\newtheorem{theorem}{Theorem}[section]
\newtheorem{lemma}[theorem]{Lemma}
\newtheorem{proposition}[theorem]{Proposition}
\newtheorem{corollary}[theorem]{Corollary}
\theoremstyle{definition}
\newtheorem{definition}[theorem]{Definition}
\theoremstyle{remark}
\newtheorem{remark}[theorem]{Remark}
\raggedbottom

\title{A hereditarily atomic domain without ACCP}
\author{AI-authored manuscript}
\date{November 17, 2026}

\begin{document}
\maketitle

\begin{abstract}
    We construct a domain $D$ with fraction field $\mathbb{Q}(t)$ in which every unital subring is atomic, although $(t^2)$ starts an infinite strictly ascending chain of principal ideals.
    The element $t$ is an atom, and the successive ratios in the chain are multiplicatively independent modulo $\mathbb{Q}^{\times}$.
    The domain is almost Dedekind, every nonzero proper ideal is a finite product of radical ideals, and its Picard group is a countable direct sum of cyclic groups of order three.
    A subring of $D$ fails bounded factorization precisely when it has finite additive index.
    These subrings are exactly $\mathbb{Z} + nD$ for $n \ge 1$, and each fails the ascending chain condition on principal ideals.
    The proof constructs two $p$-adic evaluations at every rational prime and shows that a subring with unbounded decomposition lengths contains $nD$ as soon as it contains $nt$.
\end{abstract}

\section{Introduction}
\label{sec:intro}

An integral domain is \emph{atomic} if each nonzero nonunit is a product of irreducible elements, also called \emph{atoms}.
The ascending chain condition on principal ideals (ACCP) implies atomicity.
Grams constructed an atomic domain that fails ACCP \cite[Theorem~1.3]{Grams}.
Coykendall, Gotti and Hasenauer asked whether this can happen when \emph{every} subring is atomic.
They conjectured that hereditary atomicity is equivalent to hereditary ACCP \cite[Conjecture~6.1]{CGH}.
Since hereditary atomicity passes to subrings, their conjecture is equivalent to the assertion that every hereditarily atomic domain itself satisfies ACCP.
All rings in this paper are commutative domains unless indicated otherwise, and subrings contain the same identity.

\begin{maintheorem}
    \label{thm:main}
    There is a hereditarily atomic domain $D$ with fraction field $\mathbb{Q}(t)$ and unit group $\{\pm 1\}$ for which $t$ is an atom and $(t^2)$ starts an infinite strictly ascending chain of principal ideals.
    Its successive ratios can be chosen multiplicatively independent in $\mathbb{Q}(t)^{\times}/\mathbb{Q}^{\times}$.
    The domain $D$ is an almost Dedekind SP-domain, and
    \[
        \operatorname{Pic}(D) \cong \bigoplus_{j \ge 0} \mathbb{Z}/3\mathbb{Z}.
    \]
\end{maintheorem}

An \emph{almost Dedekind domain} is a domain whose localizations at maximal ideals are discrete valuation rings.
An \emph{SP-domain} is one in which every nonzero proper ideal is a finite product of radical ideals.
A domain is \emph{Pr\"ufer} if each nonzero finitely generated ideal is invertible, or equivalently if its localizations at maximal ideals are valuation rings.
Here $\operatorname{Pic}(D)$ is the group of invertible fractional ideals modulo nonzero principal fractional ideals.
The exponent three in Theorem~\ref{thm:main} is the smallest possible finite Picard exponent for an atomic Pr\"ufer domain failing ACCP, as we show in Section~\ref{sec:sharpness}.

A \emph{bounded factorization domain} (BFD) is an atomic domain in which the lengths of atom factorizations of each fixed nonzero element are bounded.
Equivalently, for each such element the lengths of all its decompositions into nonunits are bounded.
Indeed, in an atomic domain each decomposition refines to an atom factorization, and under a bound a decomposition of maximal length already consists of atoms.
An infinite strictly ascending chain of principal ideals with nonzero first term $(x)$ gives arbitrarily long decompositions of $x$ by taking successive quotients of generators.
Thus bounded factorization implies ACCP.

\begin{maintheorem}
    \label{thm:subrings}
    For the domain $D$ in Theorem~\ref{thm:main} and any unital subring $T \subseteq D$, the following conditions are equivalent.
    \begin{enumerate}
        \item $T$ fails ACCP.
        \item $T$ fails bounded factorization.
        \item The additive index $[D:T]$ is finite.
        \item $T = \mathbb{Z} + nD$ for an integer $n \ge 1$.
    \end{enumerate}
    In the last case $n = [D:T]$ and $nD$ is the conductor $(T:D)$, the largest ideal of $D$ contained in $T$.
    Every infinite-index subring of $D$ is hereditarily a BFD.
\end{maintheorem}

The construction uses the ordinary order of vanishing $v_E$ at $t = 0$ and two evaluations in $\mathbb{Q}_p$ for each rational prime $p$.
One evaluation at each prime records the divisibility of the successive ratios of the chain.
The other forces a rational scalar that occurs in a divisor supported on the first evaluations to be a sign.
The ratios have disjoint supports, while the primary values of any fixed rational function eventually equal its order at $E$.
This gives an infinite chain below an element of order two.
At order one, a congruence modulo three between adjacent coordinates forces every element to have only finitely many divisors.

The difficulty is to factor an element in a subring $T$ that need not contain any of the auxiliary rational functions used to build $D$.
From arbitrarily long decompositions of a single element of $T$ we extract divisors whose primary supports escape every finite set.
The polynomials of Section~\ref{sec:construction} ensure that these divisors first recover the full rational algebra $\mathbb{Q}D$, and then determine the conductor at each rational prime.
Theorem~\ref{thm:conductor} shows that for such a subring, every positive integer $n$ with $nt \in T$ satisfies $nD \subseteq T$.
A sequence of order-one atoms approaches $1$ modulo each such $nD$.
These atoms can be divided out inside $T$, and induction on $v_E$ gives hereditary atomicity.

\section{The domain and its divisor profiles}
\label{sec:profiles}

Write $K = \mathbb{Q}(t)$ and $v_E = \operatorname{ord}_{t = 0}$.
Enumerate the rational primes in increasing order as $p_0, p_1, \ldots$.
We use $\mathbb{Z}_p = \varprojlim_h \mathbb{Z}/p^h\mathbb{Z}$ and its fraction field $\mathbb{Q}_p$, with $\nu_p(p) = 1$.
A condition $\nu_p(x - y) \ge h$ means $x - y \in p^h\mathbb{Z}_p$, and for integral $x, y$ this is agreement of their residues modulo $p^h$.
For each $p$ we will choose two elements $\alpha_p, \beta_p \in \mathbb{Z}_p$ transcendental over $\mathbb{Q}$.
We require $\beta_p$ to be a unit and impose
\[
    \nu_{p_j}(\alpha_{p_j}) =
    \begin{cases}
        2 & j = 0,\\
        1 & j \ge 1.
    \end{cases}
\]

For every $\gamma \in \{\alpha_p, \beta_p\}$, transcendence ensures that every nonzero polynomial in $\mathbb{Q}[t]$ has nonzero value at $\gamma$.
Thus evaluation at $\gamma$ embeds $K$ into $\mathbb{Q}_p$.
It defines the discrete valuation $f \mapsto \nu_p(f(\gamma))$ on $K$ and valuation ring $V_{p, \gamma}$, with uniformizer $p$.
Write $v_j(f) = \nu_{p_j}(f(\alpha_{p_j}))$ for the \emph{primary} valuations and call the evaluations at $\beta_p$ the \emph{guards}.
The $j$th primary and its guard lie over the same prime $p_j$.
The consecutive primary indices $2n$ and $2n + 1$ lie over different primes.

The other choices are distinct monic irreducible polynomials $F_i \in \mathbb{Q}[t]$, indexed by $i \ge 1$, with $F_i(0) \ne 0$.
Set
\begin{equation}
    \label{eq:ring}
    B = \mathbb{Q}[t, F_i^{-1} \ (i \ge 1)],
    \qquad
    D = B \cap \bigcap_p \bigl(V_{p, \alpha_p} \cap V_{p, \beta_p}\bigr)
    \quad \text{inside } K.
\end{equation}
The \emph{retained horizontal points} are represented by the monic irreducible polynomials of $\mathbb{Q}[t]$ other than the $F_i$, including $t$ itself, which represents $E$.
Over $\overline{\mathbb{Q}}$, these give all finite points except the roots of the $F_i$.
The value at the point represented by $h$ is the exponent of $h$ in the numerator minus its exponent in the denominator.
Membership in $D$ means that the values at all these horizontal points and both sorts of vertical evaluations are nonnegative.
For $f \in K^{\times}$ we call their list the \emph{profile} of $f$.
In particular, $x$ divides $y$ in $D$ precisely when the profile of $x$ is coordinatewise at most that of $y$, for nonzero $x, y \in D$.

Let $e_j$ denote the sequence supported at the primary index $j$ with value one there.
Introduce elements $b_n, q_n$ of $B^{\times}$ in the order
\[
    b_0,\ b_1,\ q_0,\ b_2,\ q_1,\ b_3,\ q_2,\ \ldots.
\]
At stage $i$ the corresponding element is $F_i$ divided by a power of the product of the earlier $F_l$.
This change of generators is triangular, so the $b_n$ and $q_n$ freely generate $B^{\times}/\mathbb{Q}^{\times}$.
All these elements must have guard value zero.
Their required primary profiles, with $v(f) = (v_j(f))_{j \ge 0}$, are
\begin{equation}
    \label{eq:profiles}
    v(b_n) = 3e_{2n},
    \qquad
    v(q_n) = 2e_{2n} + e_{2n + 1} + 2e_{2n + 2}.
\end{equation}

Two calculations explain these prescriptions.
First, the supports of $v(q_{2j})$ are pairwise disjoint, and their entries are at most two.
It follows that every finite product of distinct $q_{2j}$ divides $t^2$ in $D$.
Second, put
\[
    a_0 = t,
    \qquad
    a_{n + 1} = \frac{a_n b_{n + 1}}{q_n}.
\]
Each $a_n$ has horizontal order one at $E$, all its other horizontal and guard orders vanish, and its primary profile is
\begin{equation}
    \label{eq:atoms-profile}
    v_j(a_n) =
    \begin{cases}
        0 & j < 2n,\\
        2 & j = 2n,\\
        1 & j > 2n.
    \end{cases}
\end{equation}
Indeed, $a_0$ has this profile by the choice of valuations of the $\alpha_p$.
In the recursion the entries at $2n$ and $2n + 1$ lose two and one, respectively, and the entry at $2n + 2$ gains $3 - 2 = 1$.
Thus $a_n \in D$.
We will arrange in addition that for each fixed integer $m \ge 1$, the elements $a_n$ are eventually congruent to $1$ modulo $mD$.
This condition will make them available as divisors in every subring containing $mD$.

\section{The simultaneous polynomial construction}
\label{sec:construction}

Write $U_i$ for the $i$th term of the sequence $b_0, b_1, q_0, b_2, q_1, \ldots$, beginning with $i = 1$.
At stage $i$ we shall construct $U_i = F_i / B_i$, where $F_i \in \mathbb{Q}[t]$ is monic irreducible and $B_i$ is a power of $F_1 \cdots F_{i - 1}$.
For the first stage take $B_1 = 1$.
The issue is that the subring $T$ to be studied later is arbitrary.
Lemma~\ref{lem:extraction} will produce elements of $T$ which are products of powers of the $U_h$.
The exponents are bounded in terms of one element of $T$, and the exponent of the last generator occurring is positive.
We can therefore prescribe the behaviour of all possible products in advance, using at stage $i$ the finite list
\[
    \mathcal{W}_i
    = \left\{U_i^{a} \prod_{h < i} U_h^{k_h}
      \ \middle|\ a, k_h \in \mathbb{Z},\ 1 \le a \le i,\ |k_h| \le i\right\}.
\]
Here and below a retained geometric point is a point of $\overline{\mathbb{Q}}$ at which every $F_i$ is nonzero.

The pole and separation conditions below will recover $\mathbb{Q}T$ in Section~\ref{sec:curve}.
The two arithmetic conditions have different jobs in Section~\ref{sec:conductor}.
Coefficient congruences put a power of $(t - \alpha_p)(t - \beta_p)$ in the local conductor, and large derivatives remove repeated factors.
All derivatives are formal derivatives with respect to $t$.

\begin{proposition}[Simultaneous realization]
    \label{prop:construction}
    The polynomials $F_i$, the powers $B_i$, and the arguments $\alpha_p, \beta_p$ can be chosen with all of the following properties.
    \begin{enumerate}
        \item The $U_i$ have the primary valuations in \eqref{eq:profiles} and valuation zero at every guard.
        For each nonzero $f \in \mathbb{Q}(t)$, its primary valuations eventually equal $v_E(f)$, and its guard valuations vanish with finitely many exceptions.
        The elements $a_0 = t$ and $a_{n + 1} = a_n b_{n + 1} / q_n$ belong to $D$ and satisfy $a_n - 1 \in mD$ eventually for each integer $m \ge 1$.
        \item Each $w \in \mathcal{W}_i$ has the form
        \begin{equation}
            \label{eq:scheduled-word}
            w = F_i^{a} / H_{i, w},
        \end{equation}
        where $H_{i, w}$ is a monic product of earlier polynomials, fixed before $F_i$ is chosen.
        Each $w$ has a pole at infinity and at every root of every earlier $F_h$.
        \item For each pair of distinct retained geometric points $\lambda, \mu$, all $w \in \mathcal{W}_i$ satisfy $w(\lambda) \ne w(\mu)$ once $i$ is sufficiently large.
        For each retained $\lambda$, eventually every such $w$ satisfies $w'(\lambda) \ne 0$.
        At either argument $\gamma$ over any fixed prime $p$,
        \[
            \max_{w \in \mathcal{W}_i} \nu_p(w'(\gamma)) \longrightarrow -\infty.
        \]
        \item Fix a prime $p$ and an integer $s \ge 0$.
        Set $A_s = \mathbb{Z}_p[p^s t]$, $L_p(t) = (t - \alpha_p)(t - \beta_p)$ and $g = p^{2s} L_p(t) \in A_s$.
        For every sufficiently large $i$ and every listed word $w = F_i^a/H_{i, w}$,
        \begin{equation}
            \label{eq:transfer-module}
            g^a \in p^i (A_s + A_s w).
        \end{equation}
        Thus $g^a/p^i$ is a sum $u + v w$ with $u, v \in A_s$.
        This formulation will allow us to use the estimate in any subring containing $A_s$ and $w$.
    \end{enumerate}
\end{proposition}

\begin{proof}
    \emph{The denominator and degree.}
    Suppose the preceding stages are complete.
    Express each of the finitely many products $\prod_{h < i} U_h^{k_h}$ on our list as a monomial in the preceding $F_h$ with integer exponents.
    Take $B_i = (F_1 \cdots F_{i - 1})^{K_i}$, with $K_i$ larger than the absolute values of all these exponents.
    If the older product has exponent $m_h$ on $F_h$, then $H_{i, w}$ has exponent $aK_i - m_h > 0$ on it.
    For example, at stage two we may take $K_2 = 3$.
    The listed words are $U_2^a U_1^k = F_2^a/F_1^{3a - k}$ with $1 \le a \le 2$ and $|k| \le 2$.
    Even the largest positive older exponent leaves a factor of $F_1$ in the denominator.
    Choose an integer $d_i \ge 5$ greater than all preceding degrees and large enough that $a d_i > \deg H_{i, w}$ throughout the list.
    Once we have produced an irreducible $F_i$ of degree $d_i$, it will be distinct from every predecessor, and these inequalities will give all the asserted poles.

    \emph{Assignment of the primes.}
    At the start of stage $i$, the earlier evaluation arguments are fixed.
    At every still unassigned prime the earlier polynomials have integral coefficients and unit constant terms.
    This is vacuous at stage one and will follow at each later stage from the choice of coefficient ring and constant term in \eqref{eq:raw-constant}.
    Thus every earlier polynomial is already a unit at any new primary argument in $p\mathbb{Z}_p$.
    For a new guard we must also avoid the earlier roots modulo $p$.

    Take an increasing sequence of integer cutoffs, choosing $M_i > 2^{4 + \sum_{h \le i} d_h}$ at this stage.
    Let $S_i$ consist of the primes up to $M_i$, the previously assigned primes, and the primes where the desired primary profile of $U_i$ is positive.
    For each newly assigned prime choose a transcendental element $\alpha_p \in \mathbb{Q}_p$ of the required valuation.
    Every such valuation class is uncountable, whereas the elements algebraic over $\mathbb{Q}$ are countable.
    At the first stage the finitely many primes in $S_1$ may have any nonzero guard residue.
    At each later stage, a newly assigned $p$ exceeds $M_{i - 1}$.
    Set $u = \lfloor\log_2 p\rfloor$ and choose a guard residue represented by an integer $w_p \in [u, 2u]$ at which no earlier $F_h$ vanishes modulo $p$.
    Since $p > M_{i - 1}$, this interval supplies more distinct nonzero residues than the sum of the preceding degrees.
    Choose a transcendental lift $\beta_p \in \mathbb{Z}_p$ of the selected residue.
    Both arguments are now fixed permanently.

    Choose a fresh prime $\ell > M_i$ outside $S_i$, and assign its arguments by the same rule, requiring also $w_\ell \not\equiv 1 \pmod{\ell}$.
    The cutoff leaves enough residues for this additional exclusion, including at the first stage.
    We use $\ell$ to enforce irreducibility at this stage and include it among the previously assigned primes at every later stage.
    We choose the coefficients in $\mathbb{Z}[p^{-1} : p \in S_i]$.
    The constant term $C_i = F_i(0)$ will have the form
    \begin{equation}
        \label{eq:raw-constant}
        C_i = (-1)^{d_i} \prod_{p \in S_i} p^{-N_{p, i}},
        \qquad N_{p, i} \in (\ell - 1)\mathbb{Z}_{> 0}.
    \end{equation}
    We determine the exponents $N_{p, i}$ after fixing the local targets for the two evaluations.
    Whatever these exponents are, $C_i$ is a unit outside $S_i$ and $C_i \equiv (-1)^{d_i} \pmod{\ell}$.
    At $\ell$ we impose
    \[
        F_i(t) \equiv (t - 1)^{d_i} \pmod{\ell},
        \qquad F_i(1) \equiv \ell \pmod{\ell^2}.
    \]
    The first condition is compatible with the chosen constant term.
    After the coefficients are specified modulo $\ell$, changing the linear coefficient by a multiple of $\ell$ achieves the second condition.
    These conditions make $F_i(t + 1)$ Eisenstein at $\ell$.
    They also make $F_i$ a unit at both arguments over $\ell$, since their residues are distinct from $1$.
    Its denominator $B_i$ is a unit there by our choices of arguments.

    \emph{Interpolation at the assigned arguments.}
    At each $p \in S_i$ and each $\gamma \in \{\alpha_p, \beta_p\}$, let $h$ be the required valuation of $U_i(\gamma)$, taking $h = 0$ at a guard.
    Prescribe a local target $\eta_{p, \gamma} = B_i(\gamma) p^h$ for $F_i(\gamma)$, with the following refinement when $U_i = q_n$.
    At a branch where $h = 0$ and $a_n(\gamma)b_{n + 1}(\gamma)$ is a unit, use instead
    \[
        \eta_{p, \gamma} = B_i(\gamma) a_n(\gamma) b_{n + 1}(\gamma).
    \]
    The values on the right have already been defined at that stage.
    Every target is nonzero, since evaluation at either transcendental argument is injective on $\mathbb{Q}(t)$.
    Require $F_i(\gamma)/\eta_{p, \gamma} \in 1 + p^m\mathbb{Z}_p$, taking $m = 1$ for an ordinary target and $m = \max(1, n)$ for a refined one.
    This preserves the required valuation.
    At a refined target $a_{n + 1}(\gamma) = \eta_{p, \gamma}/F_i(\gamma)$, so the same relative congruence gives $\nu_p(a_{n + 1}(\gamma) - 1) \ge n$.

    For each $p \in S_i$ let $A_p \in \mathbb{Q}_p[t]$ be the unique polynomial of degree at most one with $A_p(\gamma) = \eta_{p, \gamma} - \gamma^{d_i}$ at the two arguments.
    Then $t^{d_i} + A_p(t)$ has the prescribed values.
    Adding a scalar multiple of $L_p$ leaves both values fixed and changes the derivative at each of its two simple roots.
    To prescribe the constant term as well, use $L_p(0) \ne 0$ and put
    \begin{align*}
        J_{p, i}(t) &= t^{d_i} + A_p(t) - A_p(0)L_p(t)/L_p(0), \\
        F_{p, i}(t) &= C_i L_p(t)/L_p(0) + J_{p, i}(t).
    \end{align*}
    Here $J_{p, i}$, the targets $F_{p, i}(\gamma) = \eta_{p, \gamma}$, and all the $H_{i, w}$ are fixed independently of the exponents in \eqref{eq:raw-constant}.
    For each $p \in S_i$ choose an integer $h_p \ge 0$ such that multiplying each polynomial in the finite list
    \[
        J_{p, i}(Z/p^s), \qquad H_{i, w}(Z/p^s)
        \quad (0 \le s \le i,\ w \in \mathcal{W}_i)
    \]
    by $p^{h_p}$ gives coefficients in $\mathbb{Z}_p$.
    Take $N_{p, i} \ge h_p + i$.
    Here $N_{p, i} = -\nu_p(C_i)$, so division by $C_i$ raises each coefficient valuation by $N_{p, i}$.
    Since $N_{p, i} - h_p \ge i$, this puts both $J_{p, i}(Z/p^s)/C_i$ and every $H_{i, w}(Z/p^s)/C_i$ in $p^i\mathbb{Z}_p[Z]$.

    We may increase these exponents further before making the rational coefficient approximation.
    At either $\gamma$ over $p$ the derivative of $F_{p, i}^{a}/H_{i, w}$ is affine in $C_i$.
    The coefficient of $C_i$ in this derivative is
    \[
        \frac{a \eta_{p, \gamma}^{a - 1} L_p'(\gamma)}
             {L_p(0) H_{i, w}(\gamma)}.
    \]
    It is nonzero because the roots of $L_p$ are simple and every factor of the denominator is nonzero at $\gamma$.
    For each derivative on the finite list there is therefore a lower bound on $-\nu_p(C_i)$ beyond which the linear term has smaller valuation than the constant term and has valuation less than $-i$.
    Choose the $N_{p, i}$ to exceed all the required local bounds and to be multiples of $\ell - 1$.
    These choices are independent across $p$, since the other factors in \eqref{eq:raw-constant} are $p$-adic units.

    \emph{Rational coefficients and geometric conditions.}
    At each $p \in S_i$ require the free coefficients of $F_i$ to be sufficiently close to those of $F_{p, i}$ to preserve all the valuation, return, and derivative conditions just imposed.
    Also require for every $0 \le s \le i$ that
    \[
        \bigl(F_i(Z/p^s) - F_{p, i}(Z/p^s)\bigr) / C_i
            \in p^i \mathbb{Z}_p[Z].
    \]
    These are finitely many open conditions and contain a product of coefficient balls.
    These local conditions can be met together with the conditions at $\ell$ by coefficients in $\mathbb{Z}[p^{-1} : p \in S_i]$.
    Indeed, first multiply the local coefficient targets by a sufficiently large common product of powers of primes in $S_i$.
    The resulting conditions contain ordinary integral congruence classes at these primes and at $\ell$, so the integer Chinese remainder theorem supplies a solution after this scaling.
    Once one solution is fixed, each free coefficient may be varied independently by any integer multiple of a sufficiently divisible common integer.
    Thus the available coefficient vectors contain a translate of a full integer grid in the affine space of monic degree-$d_i$ polynomials with constant term $C_i$.
    Such a grid is Zariski dense over $\overline{\mathbb{Q}}$, by successive application of the univariate root bound.

    Enumerate the elements of $\overline{\mathbb{Q}}$ once and for all.
    Consider the first $i$ of them, discarding any root of a preceding $F_h$, and require $F_i$ to be nonzero at each remaining point.
    At the same stage require every $w$ on our list to distinguish each pair of these points and to have nonzero derivative at each of them.
    We check that each forbidden equality cuts out a proper algebraic subset of the free coefficient space.
    Nonvanishing of $F_i(\lambda)$ is automatic at zero and can be arranged at nonzero $\lambda$ by varying the linear coefficient.
    If $\lambda, \mu$ are distinct and nonzero, their evaluation rows on the coefficients of $t, t^2$ are $(\lambda, \lambda^2)$ and $(\mu, \mu^2)$, which are independent.
    Thus $F_i(\lambda)/F_i(\mu)$ varies nontrivially.
    If either point is zero, its value is the fixed nonzero constant and the other value varies.
    As the values of $H_{i, w}$ at these points are fixed and nonzero, equality of the two values of $F_i^a/H_{i, w}$ gives a proper exclusion.
    At nonzero $\lambda$ the rows $(\lambda, \lambda^2)$ and $(1, 2\lambda)$ are independent, so the logarithmic derivative $F_i'(\lambda)/F_i(\lambda)$ varies.
    At zero it is the free linear coefficient divided by $C_i$.
    Vanishing of $w'(\lambda)$ would require $aF_i'(\lambda)/F_i(\lambda) = H_{i, w}'(\lambda)/H_{i, w}(\lambda)$, again a proper exclusion.
    There are only finitely many exclusions at any one stage, and the grid meets their common complement.
    This completes the stage.
    A point on the list which is still retained at stage $i$ remains on all subsequent lists and is protected from all subsequent numerators.
    In particular every point retained in the final ring is tested at every sufficiently late stage, giving the geometric conclusions.

    \emph{The divisibility estimate.}
    Fix $p$ and $s$ as in part~4.
    At every sufficiently late stage we have $p \in S_i$ and $s \le i$.
    The formula for the interpolating polynomial gives
    \[
        \frac{F_{p, i}}{C_i} = \frac{L_p}{L_p(0)} + \frac{J_{p, i}}{C_i}.
    \]
    The coefficient bounds for $J_{p, i}$ and $H_{i, w}$ and the rational approximation to $F_{p, i}$ therefore give
    \begin{equation}
        \label{eq:local-coefficient-estimate}
        \frac{F_i(Z/p^s)}{C_i} - \frac{L_p(Z/p^s)}{L_p(0)}
        \in p^i\mathbb{Z}_p[Z],
        \qquad
        \frac{H_{i, w}(Z/p^s)}{C_i} \in p^i\mathbb{Z}_p[Z].
    \end{equation}
    Put $\lambda = p^{2s}L_p(0) \in \mathbb{Z}_p$ and substitute $Z = p^s t$.
    We have $g = (p^s t - p^s\alpha_p)(p^s t - p^s\beta_p) \in A_s$.
    The first estimate shows that $f_i = \lambda F_i(t)/C_i$ satisfies $f_i - g \in p^i A_s$, hence $f_i^a - g^a \in p^i A_s$.
    On the other hand, clearing the denominator of the word gives
    \[
        f_i^a = w\,\lambda^a C_i^{1 - a}\,\frac{H_{i, w}(t)}{C_i}
        \in p^i A_s w .
    \]
    Indeed $a \ge 1$, $1/C_i \in \mathbb{Z}_p$, and the second estimate in \eqref{eq:local-coefficient-estimate} applies to its last factor.
    Subtraction gives \eqref{eq:transfer-module} at this single finite stage.

    \emph{Passage to all primes.}
    At each prime assigned after stage $i$, the coefficients of $F_i$ are integral and its constant term is a unit.
    Thus $F_i$ takes a unit value at the primary argument there.
    Our residue exclusions give a unit at its guard as well.
    Consequently the prescribed valuation vectors for $U_i$ hold even at primes whose arguments had yet to be chosen.
    For a fixed nonzero polynomial $V \in \mathbb{Z}[t]$, the guard representatives $w_p$ tend to infinity and are $O(\log p)$ apart from the finite initial batch.
    It follows that $0 < |V(w_p)| < p$ for all sufficiently large $p$.
    Thus $V(\beta_p)$ is a unit at these primes.
    Applying this to numerators and denominators proves the guard assertion for every rational function.
    At primaries write a fixed such function as $t^k V(t)/W(t)$ with $V, W \in \mathbb{Z}[t]$ and $V(0)W(0) \ne 0$.
    At every sufficiently late prime these two constant terms are units and $\nu_p(\alpha_p) = 1$, giving primary valuation $k$.

    Finally, the profiles and the recursion give $a_n$ primary values zero at indices $j < 2n$, two at $j = 2n$, and one at $j > 2n$.
    Its guard valuations are zero and $a_n \in B$, so $a_n \in D$.
    Fix $m \ge 1$.
    For each $p \mid m$, write $p = p_j$.
    At every sufficiently late $q_n$ stage the prime has been assigned and $j < 2n$.
    Formula~\eqref{eq:atoms-profile} gives primary value zero for $a_n$ there, and \eqref{eq:profiles} gives zero for $b_{n + 1}$ and $q_n$ there.
    All three are also units at the guard, so both branches receive the refined target.
    Taking $n \ge \nu_p(m)$ gives integral values of $(a_{n + 1} - 1)/m$ at both branches.
    At all other primes $m$ is a unit, and the same quotient belongs to $B$ since $m$ is a nonzero rational constant.
    The intersection defining $D$ gives $(a_{n + 1} - 1)/m \in D$ as asserted.
\end{proof}

\section{Maximal ideals, the Picard group, and a chain}
\label{sec:structure}

Write $P_j$ for the centre in $D$ of the primary evaluation over $p_j$, and $G_p$ for the centre of the guard over $p$.
For a monic irreducible $h \in \mathbb{Q}[t]$ which is not inverted in $B$, write $H_h = D \cap hB$.
Thus $H_t$ is the centre of $E$.
We call the $P_j$ and $G_p$ vertical points and the $H_h$ horizontal points.
For each of these points we use the normalized valuation on $\mathbb{Q}(t)$.

\begin{lemma}
    \label{lem:stalks}
    The fraction field of $D$ is $\mathbb{Q}(t)$, its rational constants are $\mathbb{Z}$, and inverting the nonzero integers in $D$ gives $B$.
    Its maximal ideals are precisely $P_j$, $G_p$, and $H_h$.
    Their localizations are the corresponding discrete valuation rings.
    In particular, $D$ is an almost Dedekind domain.
\end{lemma}

\begin{proof}
    All arguments of evaluation are integral at their respective primes, so $\mathbb{Z}[t] \subseteq D$.
    A rational constant belongs to $D$ precisely when it is integral at every rational prime.
    For $x \in B \setminus \{0\}$, its order at $E$ is nonnegative.
    By Proposition~\ref{prop:construction}, its primary values eventually equal that order, and its guard values eventually vanish.
    A positive integer clears the finitely many negative vertical values of $x$.
    This proves the assertions about the fraction field, constants, and localization at the integers.

    Set
    \[
        f_{2n} = b_n,
        \qquad
        f_{2n + 1} = \frac{q_n^3}{b_n^2 b_{n + 1}^2}.
    \]
    By \eqref{eq:profiles}, $f_j$ has value three at $P_j$ and zero at every other displayed point.
    The function $p_j^3/f_j$ has value three at $G_{p_j}$ and zero at every other displayed point.
    In particular both functions belong to $D$.
    If $V$ is either point over $p$ and $f$ is its corresponding function, then $\sqrt{fD} = V$.
    Indeed, for $x \in D$ with $v_V(x) > 0$ we have $x^3/f \in D$, and the reverse inclusion follows by applying $v_V$.
    Reduction at the argument maps $D$ onto $\mathbb{F}_p$, since $D$ contains $\mathbb{Z}$.
    Its kernel $V$ is maximal.
    Every prime of $D$ containing $p_j$ contains one of the factors in $p_j^3 = f_j(p_j^3/f_j)$, so it is one of the two displayed points over $p_j$.

    The valuation ring at $V$ contains $D_V$.
    Conversely, suppose $a, b \in D$, $b \ne 0$, and $v_V(a/b) \ge 0$.
    The element $s = b^3/f^{v_V(b)}$ belongs to $D \setminus V$.
    At each other displayed point, $f$ has value zero and $s(a/b) = ab^2/f^{v_V(b)}$ is integral.
    It is integral at $V$ as well, so $a/b \in D_V$.
    This valuation ring is discrete with value group $\mathbb{Z}$ because $v_V(p) = 1$.

    The primes disjoint from the nonzero integers correspond to the primes of the localization $B$.
    The latter is a localization of the PID $\mathbb{Q}[t]$.
    For each nonzero surviving prime $hB$, its contraction $H_h$ is maximal among the primes disjoint from the nonzero integers.
    To check maximality in $D$, choose a positive integer $d$ such that $g = dh$ belongs to $D$.
    At any vertical point $V$, the associated function $f$ above is a unit in $B$.
    Thus $g^3/f^{v_V(g)}$ belongs to $D \cap hB$ and has value zero at $V$.
    Hence $H_h$ lies in no vertical maximal ideal.
    Its localization is $B_{hB} = \mathbb{Q}[t]_{(h)}$.
\end{proof}

For an invertible fractional ideal $J$, let $v_M(J)$ be the exponent of $J D_M$ in the DVR $D_M$.
We call the tuple of these exponents its profile.
Write $e_M$ for the tuple with value one at $M$ and zero elsewhere, abbreviating $e_{P_j}$ to $e_j$.
Let $\lambda$ have value one at $E$ and at every primary, and zero at all other points.

\begin{proposition}
    \label{prop:ideals}
    The profiles of the invertible fractional ideals are precisely the integer tuples whose primary entries eventually equal their entry at $E$ and whose entries at the guards and the other horizontals have finite support.
    The domain $D$ is an SP-domain, that is, every nonzero proper ideal is a finite product of radical ideals.
    Moreover,
    \[
        D^\times = \{1, -1\},
        \qquad
        \operatorname{Pic}(D) \cong \bigoplus_{n \ge 0} \mathbb{Z}/3\mathbb{Z}.
    \]
\end{proposition}

\begin{proof}
    Ideals are determined by their localizations at all maximal ideals.
    The profile map on invertible fractional ideals is therefore injective and takes products to sums.
    By Lemma~\ref{lem:stalks}, $D$ is Pr\"ufer, so every nonzero finitely generated fractional ideal is invertible.
    Its profile is the coordinatewise minimum of the profiles of a finite generating set.
    The eventual values of rational functions from Proposition~\ref{prop:construction} imply the claimed restrictions.

    At a vertical point $V$ over $p$, take the function $f$ used in Lemma~\ref{lem:stalks}.
    The ideal $(p, f)$ has profile $e_V$, since $p$ has value one there and $f$ has value zero at every other point.
    For a horizontal point $H_h$ other than $E$, the principal fractional ideal $hD$ has horizontal profile $e_{H_h}$ and finitely many nonzero vertical entries, since $h(0) \ne 0$.
    Multiplying by integer powers of the vertical singleton ideals cancels these entries and realizes $e_{H_h}$.
    Finally $tD$ has profile $\lambda + e_0$.
    These singleton profiles and $\lambda$ generate every tuple in the statement.

    To factor an arbitrary nonzero proper ideal $I$, let $m_M = v_M(I D_M)$, the least $v_M(x)$ for nonzero $x \in I$.
    The profile of any fixed nonzero element of $I$ bounds all $m_M$ by a single integer $N$.
    For each nonzero $x \in I$ and integer $k \ge 1$, the indicator of the set of points where $v_M(x) \ge k$ satisfies the restrictions just proved.
    Let $J(x, k)$ be the invertible ideal with this profile.
    It is integral because every entry is nonnegative.
    The ideal $J_k = \sum_{0 \ne x \in I} J(x, k)$ has exponent one at $M$ if $m_M \ge k$, and zero otherwise.
    Each $J_k$ is radical, as can be checked in the maximal localizations.
    At each $M$ the exponents of $J_1, \ldots, J_N$ sum to $m_M$, giving $I = J_1 \cdots J_N$.

    To compute the Picard group, quotient the group of profiles by the principal profiles.
    Subtract the profiles of powers of the surviving irreducibles $h \ne t$ to clear the horizontal entries away from $E$.
    Then subtract the profiles of powers of rational primes to clear the guards.
    Every class now has a representative supported on the primaries and $E$.
    A rational function with principal profile supported there has the form
    \[
        c t^m \prod_{n \ge 0} b_n^{s_n} q_n^{u_n},
        \qquad
        c \in \mathbb{Q}^\times,\quad m, s_n, u_n \in \mathbb{Z},
    \]
    with only finitely many nonzero $s_n, u_n$.
    Indeed, factorization in $\mathbb{Q}[t]$ allows only $t$ and the inverted $F_i$ to occur.
    Replacing each $F_i$ by the associated generator times its denominator in earlier $F_\ell$ gives the displayed form, working backwards from the largest index.
    The functions $t$, $b_n$, and $q_n$ have value zero at every guard, so $\nu_p(c) = 0$ for every $p$, giving $c = \pm 1$.
    If the function is a unit of $D$, its $E$-value gives $m = 0$, its odd primary values then give every $u_n = 0$, and its even primary values give every $s_n = 0$.

    The principal profile of $t$ eliminates the entry at $E$, leaving a finitely supported primary tuple.
    The profile of $q_n$ eliminates the odd entry $2n + 1$ with coefficient one and does not change any other odd entry.
    The remaining classes are generated by $e_{2n}$.
    In any principal relation supported on the even entries, the same $E$ and odd tests force $m = 0$ and every $u_n = 0$.
    The remaining relations come exactly from the $b_n$, with profiles $3e_{2n}$.
    This gives the asserted direct sum.
\end{proof}

\begin{corollary}
    \label{cor:chain}
    In $D$, the ideal $t^2 D$ starts an infinite strictly ascending chain of principal ideals.
    Its successive ratios are multiplicatively independent modulo $\mathbb{Q}^\times$.
\end{corollary}

\begin{proof}
    For each integer $j \ge 0$, put
    \[
        x_j = \frac{t^2}{\prod_{k = 0}^{j - 1} q_{2k}},
        \qquad x_0 = t^2.
    \]
    The primary supports of $q_{2k}$ are the disjoint triples $\{4k, 4k + 1, 4k + 2\}$, and their largest entry is two.
    All primary values of $t^2$ are at least two.
    The denominators are units at every other displayed point, so $x_j \in D$.
    Since $x_j = q_{2j} x_{j + 1}$ and $q_{2j}$ is a nonunit, $x_jD \subsetneq x_{j + 1}D$.
    If a finite product of integral powers of the $q_{2k}$ is a rational constant $c$, evaluating the guards gives $c = \pm 1$.
    Evaluating its primary coordinate $4k + 1$ now shows that the exponent of $q_{2k}$ is zero for every $k$.
\end{proof}

\section{Recovering the rational curve from unbounded lengths}
\label{sec:curve}

Fix a unital subring $T \subseteq D$ and a nonzero element $x \in T$.
A decomposition of length $m$ means an equality $x = y_1 \cdots y_m$ with each $y_j$ a nonunit of $T$.
The factors need not be atoms.
By Proposition~\ref{prop:ideals}, $D^{\times} = \{1, -1\}$, so every nonunit of $T$ remains a nonunit of $D$.
In particular, some defining valuation of $D$ is positive on each such factor.

\begin{lemma}
    \label{lem:extraction}
    Suppose that $x$ has decompositions of arbitrarily large length in $T$, and set $r = \sup_{j \ge 0} v_j(x)$.
    For every finite initial segment of the primary indices there is a divisor $y$ of $x$ in $T$ whose valuation is positive on a nonempty finite set of primaries beyond that segment and zero at all other defining places.
    Up to sign, $y$ is a word in the $b_n$ and $q_n$ with integer exponents of absolute value at most $2r$ and positive exponent on the last generator occurring in construction order.
    Such divisors can be chosen with their last stages tending to infinity.
\end{lemma}

\begin{proof}
    The primary values of $x$ are eventually $v_E(x)$, so $r$ is finite.
    Form a finite set $\mathcal{V}$ of places consisting of $E$, the specified primary segment, and the guards and other retained horizontal places at which $x$ has positive value.
    Finiteness at the guards follows from Proposition~\ref{prop:construction}.
    In any decomposition of $x$, all values of all factors are nonnegative.
    At most $\sum_{v \in \mathcal{V}} v(x)$ factors have positive value somewhere in $\mathcal{V}$, since each such factor contributes at least one to this sum.
    Take a longer decomposition and a factor $y$ undetected on $\mathcal{V}$.
    At every guard and every horizontal place outside $\mathcal{V}$ the value of $x$ is zero, hence so is the value of $y$.
    Since $v_E(y) = 0$, its primary values are eventually zero as well.
    The factor $y$ is a nonunit, so its remaining primary support is nonempty.

    Vanishing at every retained horizontal place makes $y$ a unit of $B$.
    Unique factorization in $\mathbb{Q}[t]$ writes a unit of $B$ uniquely as a rational constant times a finite product of integer powers of the $F_i$.
    At each stage the normalized generator is $F_i$ divided by a product of powers of earlier $F_j$, so successive substitution gives a unique expression in the normalized generators with the same constant.
    Those generators are units at all guards, and the value of the constant at the guard over $p$ is its ordinary $p$-adic valuation.
    Thus the constant is a sign.
    Write
    \[
        y = \pm \prod_{n \ge 0} b_n^{s_n} q_n^{u_n},
        \qquad g_j = v_j(y), \qquad u_{-1} = 0.
    \]
    By the profiles in \eqref{eq:profiles},
    \[
        u_n = g_{2n + 1},
        \qquad
        3s_n = g_{2n} - 2u_n - 2u_{n - 1}.
    \]
    All $g_j$ lie between zero and $r$ because $y$ divides $x$ in $D$.
    It follows that $0 \le u_n \le r$ and $-4r \le 3s_n \le r$, which give the asserted bound.
    If the last generator is $q_n$, its nonzero exponent $u_n$ is positive.
    If it is $b_n$, both $q_n$ and, for $n \ge 1$, $q_{n - 1}$ occur later in construction order.
    Their exponents therefore vanish, and the nonzero exponent $s_n$ equals $g_{2n}/3 > 0$.
    Finally, any fixed initial list of generators meets only finitely many primary places.
    Requiring the nonempty support of $y$ to lie beyond all these places forces its last stage beyond the list.
\end{proof}

The exponents in Lemma~\ref{lem:extraction} may have either sign away from the last stage.
The predecessor powers in Proposition~\ref{prop:construction} were chosen for precisely these words.
Put $R = 2r$ and choose such a sequence of divisors with strictly increasing last stages $i \ge R$.
After multiplying by signs, the resulting elements of $T$ occur on the lists at those stages as $F_i^{a}/H$ with $1 \le a \le R$.
We shall use the poles and geometric values of this sequence here and its $p$-adic estimates in Section~\ref{sec:conductor}.

It is possible to distinguish every pair of geometric points and still miss $t$.
For example, $k[t^2, t^3] \subsetneq k[t]$ over $k = \overline{\mathbb{Q}}$ separates points, since away from zero the quotient of $t^3$ by $t^2$ recovers the coordinate.
At zero every element of the smaller ring has derivative zero, and the ideal $(t^2, t^3)$ cuts out the quotient $k[t]/(t^2)$ in the larger ring.
The next lemma uses a nonzero derivative to rule out this last multiplicity.

\begin{lemma}
    \label{lem:curve-recognition}
    Let $A$ be a $\mathbb{Q}$-subalgebra of $B$.
    Suppose that each omitted geometric point, including infinity, is a pole of some member of $A$.
    Suppose also that members of $A$ separate every pair of distinct retained geometric points and that at each retained point some member of $A$ has nonzero first derivative.
    Then $A = B$.
\end{lemma}

\begin{proof}
    Let $P/Q \in A$, with $P, Q \in \mathbb{Q}[t]$ coprime.
    Choose polynomials $U, V \in \mathbb{Q}[t]$ with $UP + VQ = 1$.
    Then $1/Q = U(P/Q) + V \in A[t]$.
    For each $F_i$, a pole at a root of $F_i$ gives a denominator $Q$ divisible by $F_i$, whence $1/F_i = (Q/F_i)/Q \in A[t]$.
    Thus $B = A[t]$.
    At infinity the pole condition gives $h = P/Q \in A$ with $\deg{P} > \deg{Q}$.
    The polynomial $P(X) - hQ(X) \in A[X]$ has $t$ as a root, and its leading coefficient is a nonzero rational number.
    Dividing through by this coefficient shows that $t$ is integral over $A$, so $B$ is a finite $A$-module.

    Put $k = \overline{\mathbb{Q}}$ and extend scalars to obtain a finite extension $A_k \subseteq B_k$, where $A_k = k \otimes_{\mathbb{Q}} A$ and $B_k = k \otimes_{\mathbb{Q}} B$.
    We have $B_k = k[t, F_i^{-1} : i \ge 1]$.
    Its maximal ideals are exactly the ideals $(t - \alpha)B_k$ at retained points $\alpha$.
    Indeed $0$ as a geometric point is retained, so $B_k$ is not a field, and its nonzero primes come from the noninverted linear factors in $k[t]$.

    Fix a maximal ideal $\mathfrak{m}$ of $A_k$.
    By lying over it has a maximal ideal $(t - \alpha)B_k$ above it.
    The residue field of $\mathfrak{m}$ embeds in $k$ and contains the constants $k$, hence equals $k$.
    Evaluations at two points above $\mathfrak{m}$ would then agree on $A_k$, so the separation hypothesis makes $\alpha$ unique.
    Choose $h \in A$ with $h'(\alpha) \ne 0$.
    The ideal $\mathfrak{m}B_k$ contains $h - h(\alpha)$, a nonzero element with a simple zero at $\alpha$.
    In the PID $B_k$ the only maximal ideal containing $\mathfrak{m}B_k$ is $(t - \alpha)B_k$, by uniqueness of the point above $\mathfrak{m}$.
    Hence $\mathfrak{m}B_k = (t - \alpha)^e B_k$ for some $e \ge 1$, and the element with a simple zero forces $e = 1$.
    Thus the fibre $B_k/\mathfrak{m}B_k$ equals $k$.

    In particular the map $A_k/\mathfrak{m} \to B_k/\mathfrak{m}B_k$ is onto.
    For the finite $A_k$-module $M = B_k/A_k$ this gives $M/\mathfrak{m}M = 0$.
    Nakayama's lemma yields $M_{\mathfrak{m}} = 0$ for every maximal ideal $\mathfrak{m}$, so $A_k = B_k$.
    Faithful flatness of $k$ over $\mathbb{Q}$ gives $A = B$.
\end{proof}

\begin{proposition}
    \label{prop:generic}
    If some nonzero element of $T$ has decompositions of arbitrarily large length in $T$, then $\mathbb{Q}T = B$.
    In particular $nt \in T$ for some positive integer $n$.
\end{proposition}

\begin{proof}
    Apply Lemma~\ref{lem:extraction} and take the listed words of increasing last stages described above.
    They have a pole at infinity, and for every fixed $j$ a sufficiently late word has a pole at every root of $F_j$.
    For each fixed pair of distinct retained points all sufficiently late words distinguish the pair, and at each fixed retained point all sufficiently late words have nonzero derivative, by Proposition~\ref{prop:construction}.
    These words belong to $T$, so Lemma~\ref{lem:curve-recognition} applies to $A = \mathbb{Q}T$.
    Since $\mathbb{Q}T$ is the localization of $T$ at the nonzero integers and contains $t$, a positive integer clears the denominator of $t$.
\end{proof}

\section{An integer in the conductor}
\label{sec:conductor}

Let $T$ be a unital subring of $D$ in which some nonzero element admits arbitrarily long decompositions into nonunits.
By Proposition~\ref{prop:generic}, $\mathbb{Q}T = B$.
In particular, some positive integer $n$ satisfies $nt \in T$.
We will prove that every such $n$ satisfies $nD \subseteq T$.

Fix a rational prime $p$ and put
\[
    T' = \mathbb{Z}_p \otimes_{\mathbb{Z}} T,
    \qquad
    D' = \mathbb{Z}_p \otimes_{\mathbb{Z}} D,
    \qquad
    B' = \mathbb{Q}_p \otimes_{\mathbb{Q}} B.
\]
The ring $B'$ is the localization of $\mathbb{Q}_p[t]$ obtained by inverting the $F_i$.
These are ordinary tensor products, so their elements are finite sums of pure tensors.
The module $\mathbb{Z}_p$ is flat over $\mathbb{Z}$ because it is torsion-free over a PID.
Tensoring $T \subseteq D \subseteq B$ therefore gives inclusions, and $\mathbb{Z}_p \otimes_{\mathbb{Z}} B = B'$ since $B$ is a $\mathbb{Q}$-algebra.
Moreover $T'[1/p] = \mathbb{Q}_p \otimes_{\mathbb{Q}} (\mathbb{Q}T) = B'$.
Neither $\alpha_p$ nor $\beta_p$ is a root of any $F_i$, since both arguments are transcendental over $\mathbb{Q}$.
Consequently evaluation at either argument extends to $B'$.
At $\gamma \in \{\alpha_p, \beta_p\}$ the value on $\sum_j u_j \otimes f_j \in D'$ is $\sum_j u_j f_j(\gamma) \in \mathbb{Z}_p$.
The scalar extension puts $\gamma$ in the coefficient field, so evaluation on $B'$ now has kernel $(t - \gamma)B'$.
The kernel of the pair of evaluations is consequently generated by $L_p(t) = (t - \alpha_p)(t - \beta_p)$.
If this entire kernel lies in $T'$, membership in $T'$ can be checked on the two values alone.

By Lemma~\ref{lem:extraction}, there are words in $T$ from arbitrarily late stages $i$ of the construction, of the form
\[
    w_i = \frac{F_i(t)^{e_i}}{H_i(t)},
    \qquad 1 \le e_i \le R,
\]
where $R = 2r$ is fixed and a possible sign has been absorbed using $-1 \in T$.
Here $r$ bounds the primary valuations of the element with unbounded decomposition lengths, and $H_i$ denotes the denominator of the particular word selected at stage $i$.

We will use the following test twice.
For $g \in B'$, the equality $B' = T'[1/p]$ gives
\[
    g B' \subseteq T'
    \quad\Longleftrightarrow\quad
    g/p^k \in T' \quad\text{for every } k \ge 0.
\]
Indeed every element of $B'$ is $u/p^k$ for some $u \in T'$.
We first apply the test to a power of $L_p$.
We then show that any repeated factor of the conductor generator could be cancelled, by taking a derivative at its root.

\begin{lemma}[Exact polynomial transfer]
    \label{lem:exact-transfer}
    One has $L_p(t)^R B' \subseteq T'$.
\end{lemma}

\begin{proof}
    Fix a positive integer $n$ with $nt \in T$ and put $s = \nu_p(n)$.
    The scalar $n/p^s$ is a unit in $\mathbb{Z}_p$, so $p^s t \in T'$ and $A_s = \mathbb{Z}_p[p^s t] \subseteq T'$.
    Put $g = p^{2s}L_p(t) \in A_s$.
    Fix any integer $k \ge 0$ and choose an extracted stage $i \ge k$ sufficiently late for \eqref{eq:transfer-module}.
    As $w_i \in T'$, that estimate gives $g^{e_i}/p^i \in A_s + A_s w_i \subseteq T'$.
    Multiplication by $p^{i - k}g^{R - e_i} \in T'$ yields $g^R/p^k \in T'$.
    The test above puts $g^R B'$ in $T'$, and $g^R B' = L_p^R B'$ since $p$ is invertible in $B'$.
\end{proof}

\begin{lemma}[The local conductor]
    \label{lem:local-conductor}
    The conductor $(T' : B') = \{f \in B' \mid fB' \subseteq T'\}$ equals $L_p(t)B'$.
\end{lemma}

\begin{proof}
    Write $I = (T' : B')$, the largest ideal of $B'$ contained in $T'$.
    Lemma~\ref{lem:exact-transfer} gives $L_p^R B' \subseteq I$.
    For $f \in I$, every $p^{-k} f$ with $k \ge 0$ belongs to $T'$ and has integral values at $\alpha_p$ and $\beta_p$.
    Both values of $f$ are therefore zero.
    The two distinct linear factors of $L_p$ survive in the localized principal ideal domain $B'$, so
    \[
        L_p^R B' \subseteq I \subseteq L_p B'.
    \]
    Thus we may write $I = h B'$ where $h$ is a product of powers of $t - \alpha_p$ and $t - \beta_p$, each with exponent between one and $R$.

    Suppose $(t - \gamma)^2$ divides $h$ at one of the two arguments.
    Put $q = h/(t - \gamma)^2$.
    We will prove $(h/(t - \gamma)) B' \subseteq T'$, contradicting the maximality of $I$.
    Choose an integer $d \ge 0$ with $p^d q \in T'$, using $T'[1/p] = B'$.
    Fix $k \ge 0$.
    Along the extracted sequence, $\nu_p(w_i'(\gamma)) \to -\infty$ by Proposition~\ref{prop:construction}.
    We may therefore choose $w_i$ with $u = 1/(p^k w_i'(\gamma)) \in p^d\mathbb{Z}_p$.
    Then $uq \in T'$, and the scalar $w_i(\gamma)$ belongs to $\mathbb{Z}_p \subseteq T'$.

    In $B'$ the first-order expansion at $\gamma$ gives
    \[
        w_i(t) - w_i(\gamma) \equiv w_i'(\gamma)(t - \gamma)
        \pmod{(t - \gamma)^2 B'}.
    \]
    To see this for a rational function, expand its numerator and denominator modulo $(t - \gamma)^2$ and invert the denominator, whose value at $\gamma$ is nonzero.
    After multiplication by $uq$ the error belongs to $h B' = I$.
    Consequently
    \[
        uq\bigl(w_i(t) - w_i(\gamma)\bigr)
        \equiv \frac{h/(t - \gamma)}{p^k} \pmod{I}.
    \]
    The left side belongs to $T'$ and $I \subseteq T'$, so $(h/(t - \gamma))/p^k \in T'$ for every $k$.
    The test above now puts $(h/(t - \gamma))B'$ in $T'$.
    This ideal strictly contains $I$ because $t - \gamma$ is a nonunit of $B'$.
    Neither factor can be repeated, so $I = L_p B'$.
\end{proof}

\begin{theorem}[The integer conductor]
    \label{thm:conductor}
    Suppose a unital subring $T \subseteq D$ has a nonzero element with unbounded decomposition lengths.
    There exists a positive integer $n$ with $nt \in T$, and every such $n$ satisfies $nD \subseteq T$.
\end{theorem}

\begin{proof}
    Existence follows from $\mathbb{Q}T = B$.
    Fix any such $n$ and any rational prime $p$.
    By Lemma~\ref{lem:local-conductor} the kernel of the value map $B' \to \mathbb{Q}_p^2$ is contained in $T'$.
    The value vectors of $1$ and $t$ form the columns of the matrix
    \[
        \begin{pmatrix}
            1 & \alpha_p \\
            1 & \beta_p
        \end{pmatrix}.
    \]
    Its determinant is a $p$-adic unit because $\alpha_p \in p\mathbb{Z}_p$ and $\beta_p \in \mathbb{Z}_p^\times$.
    Since $n, nt \in T'$, the image of $T'$ contains $n\mathbb{Z}_p^2$.
    Each element of $nD'$ has its two values in $n\mathbb{Z}_p^2$ and hence has the same values as some element of $T'$.
    Their difference lies in the kernel, so $nD' \subseteq T'$.

    For descent, take the abelian group $M = (nD + T)/T$.
    Flatness of $\mathbb{Z}_p$ over $\mathbb{Z}$ identifies $\mathbb{Z}_p \otimes_{\mathbb{Z}} M$ with $(nD' + T')/T' = 0$.
    To see that the tests at all primes detect $M$, observe that $M$ is torsion since $\mathbb{Q}T = B$.
    A nonzero element of $M$ would generate a subgroup isomorphic to $\mathbb{Z}/m\mathbb{Z}$ with $m > 1$.
    For a prime $p \mid m$, flatness would embed the nonzero group $\mathbb{Z}_p/m\mathbb{Z}_p$ in $\mathbb{Z}_p \otimes_{\mathbb{Z}} M$.
    Hence $M = 0$, which is the asserted inclusion.
\end{proof}

\section{Atomicity and the subrings of finite index}
\label{sec:orders}

Write $v_E(x)$ for the order of a nonzero $x \in D$ at $t = 0$.
Since $D^{\times} = \{\pm 1\}$, every unital subring $T$ of $D$ has the same units.
In particular, taking divisors in $D$ gives bounds on the divisors and on the lengths of decompositions in $T$.

\begin{lemma}
    \label{lem:low-order}
    Every nonzero element of $D$ of order zero or one at $E$ has only finitely many divisors up to associates.
    Each $a_n$ is an atom of $D$.
\end{lemma}

\begin{proof}
    Let $g$ be a principal valuation profile supported on finitely many primaries, and set $g_{-1} = 0$.
    The profiles of the $b_n$ and $q_n$ generate all such profiles as an abelian group.
    Each generator in \eqref{eq:profiles} satisfies
    \begin{equation}
        g_{2k - 1} + g_{2k} + g_{2k + 1} \equiv 0 \pmod{3}
        \qquad (k \ge 0).
        \label{eq:triple-congruence}
    \end{equation}
    Suppose that the entries of $g$ are zero or one for all $j \ge 2N - 1$, where $N \ge 1$.
    For $k \ge N$, the three entries in \eqref{eq:triple-congruence} must then be all zero or all one.
    Consecutive triples share an entry, and every sufficiently late triple is zero.
    It follows that $g_j = 0$ for all $j \ge 2N - 1$.
    No restriction on the signs of the entries before this boundary is needed.

    If $v_E(x) = 0$, the entire profile of $x$ is finitely supported.
    The profile of any divisor in $D$ lies between zero and the profile of $x$ at every coordinate, so there are only finitely many possibilities.
    Elements with the same profile differ by a unit.

    Suppose that $v_E(x) = 1$.
    Consider first a divisor $z$ of $x$ with $v_E(z) = 0$.
    There are only finitely many possible tuples of its values at the guards and at the retained horizontal points other than $E$, since $x$ has finite support there.
    For each tuple that occurs, fix one divisor $u$ giving that tuple.
    The profile $g$ of $z/u$ for any other divisor $z$ in the same class is supported on finitely many primaries.
    Fix $N \ge 1$ large enough that $v_j(x) = 1$ and $v_j(u) = 0$ for $j \ge 2N - 1$.
    Then $g_j = v_j(z) \in \{0, 1\}$ on this tail, so the congruence argument gives $v_j(z) = 0$ there.
    The finitely many other primary values range from zero to the corresponding values of $x$.
    This proves finiteness in each class.
    The involution $z \mapsto x/z$ exchanges divisors of order zero with divisors of order one, which proves the assertion for $x$.

    Finally, any decomposition of $a_n$ into two factors has a factor of order zero.
    The profile $g$ of such a divisor is supported on finitely many primaries, vanishes before $2n$, and is bounded by one from $2n + 1$ onward.
    Applying the triple argument with $N = n + 1$ leaves only the entry $g_{2n}$.
    Equation~\eqref{eq:triple-congruence} at $k = n$ makes this entry divisible by three, while $0 \le g_{2n} \le 2$.
    Thus the divisor is a unit, and $a_n$ is an atom.
\end{proof}

\begin{proposition}
    \label{prop:all-atomic}
    Every unital subring $T \subseteq D$ is atomic.
    If $mD \subseteq T$ for some positive integer $m$, every factorization of a nonzero nonunit $x \in T$ into atoms contains exactly $v_E(x)$ factors of order one at $E$.
    All its other factors have order zero.
\end{proposition}

\begin{proof}
    In a decomposition into nonunits, the successive partial products are pairwise nonassociate divisors of the element being decomposed.
    Consequently Lemma~\ref{lem:low-order} bounds the lengths of such decompositions in $T$ at orders zero and one.
    A decomposition of greatest length consists entirely of atoms, since splitting any of its factors increases its length.

    Suppose first that $mD \subseteq T$, and let $v_E(x) = r \ge 2$.
    We have $v_j(x) = r$ at all sufficiently late primaries.
    The primary values of $a_i$ vanish before $2i$, equal two at $2i$, and equal one afterwards.
    Its only other positive value is one at $E$.
    Hence every sufficiently late $a_i$ divides $x$ in $D$.
    By Proposition~\ref{prop:construction} we may take $i$ large enough that $a_i \equiv 1 \pmod{mD}$ as well.
    Then $a_i \in T$, and the identity
    \[
        \frac{x}{a_i} - x = \frac{x}{a_i}(1 - a_i) \in mD
    \]
    gives $x/a_i \in T$.
    The factor $a_i$ remains an atom of $T$ because both rings have the same units.
    The other factor has order $r - 1$, so induction on $r$ gives a factorization of $x$ into atoms.
    The same splitting shows that no atom of $T$ has order greater than one.
    Additivity of $v_E$ now gives the asserted count.

    For a general subring, if every nonzero nonunit has bounded lengths of decompositions into nonunits, taking a decomposition of greatest length gives atomicity.
    Otherwise Theorem~\ref{thm:conductor} gives $mD \subseteq T$ for some positive integer $m$, and the preceding argument applies.
\end{proof}

In particular $t = a_0$ is an atom.
The strict chain beginning at $(t^2)$ from Section~\ref{sec:structure} thus begins at the square of an atom.
This completes the proof of Theorem~\ref{thm:main}.

\begin{proposition}
    \label{prop:finite-index}
    For each positive integer $n$, the ring $T_n = \mathbb{Z} + nD$ is the unique unital subring of $D$ of additive index $n$.
    Its elements are precisely the $f \in D$ satisfying
    \[
        f(\alpha_p) \equiv f(\beta_p) \pmod{p^{\nu_p(n)}}
        \qquad (p \mid n).
    \]
    The additive quotient $D/T_n$ is cyclic, generated by the class of $t$, and the conductor $(T_n : D)$ equals $nD$.
\end{proposition}

\begin{proof}
    Let $p$ be a rational prime and let $h \ge 1$.
    Evaluation at $\alpha_p$ and $\beta_p$ induces an isomorphism
    \begin{equation}
        D/p^hD \longrightarrow
        (\mathbb{Z}/p^h\mathbb{Z}) \times (\mathbb{Z}/p^h\mathbb{Z}).
        \label{eq:two-residues}
    \end{equation}
    Indeed, if both evaluations of $x \in D$ are divisible by $p^h$, the quotient $x/p^h$ is integral at the two valuations over $p$.
    The denominator is a unit at every other defining valuation, so $x/p^h \in D$.
    For surjectivity, the images of $1$ and $t$ form a basis of the target module over $\mathbb{Z}/p^h\mathbb{Z}$, since $\beta_p - \alpha_p$ is a $p$-adic unit.
    The Chinese remainder theorem for the ideals $p^hD$ gives the product of these descriptions for each integer modulus $m > 1$.
    In particular $D/mD$ is finite.

    For $n > 1$ the Chinese remainder isomorphism identifies $D/nD$ with $(\mathbb{Z}/n\mathbb{Z})^2$, using the primary value in each first coordinate and the guard value in each second coordinate.
    The image of $\mathbb{Z}$ is the diagonal.
    Thus $T_n$ is the inverse image of the diagonal, which gives the stated congruences.
    Subtracting the first coordinate from the second yields a surjective additive homomorphism $D \to \mathbb{Z}/n\mathbb{Z}$ with kernel $T_n$.
    The class of $t$ maps to a unit modulo each $p \mid n$, so it generates this additive quotient and $[D:T_n] = n$.

    If a unital subring $T$ has additive index $n$, its additive quotient is killed by $n$.
    Therefore $\mathbb{Z} + nD \subseteq T$, and equality follows since both subgroups have index $n$.
    Finally the diagonal contains no nonzero ideal of $(\mathbb{Z}/n\mathbb{Z})^2$, since multiplying a diagonal element $(u, u)$ by $(1, 0)$ gives $(u, 0)$.
    Consequently $nD$ is the largest ideal of $D$ contained in $T_n$.
    For $n = 1$ these assertions give $T_1 = D$, the trivial quotient and conductor $D$.
\end{proof}

\begin{proposition}
    \label{prop:chain-height}
    Let $T \subseteq D$ have finite additive index.
    For nonzero $x \in T$, the ideal $xT$ starts an infinite strictly ascending chain of principal ideals of $T$ if and only if $v_E(x) \ge 2$.
    At every such $x$ the lengths of atomic factorizations in $T$ are unbounded.
    At orders zero and one there are only finitely many factorizations up to order and associates.
\end{proposition}

\begin{proof}
    At orders zero and one Lemma~\ref{lem:low-order} gives finitely many divisors.
    Every ideal in a chain above $xT$ is generated by a divisor of $x$.
    Also each atomic factor is a divisor, and successive partial products bound all lengths as in Proposition~\ref{prop:all-atomic}.
    This proves both finiteness assertions at these orders.

    Write $T = T_n$ by Proposition~\ref{prop:finite-index}, and suppose that $v_E(x) = r \ge 2$.
    The supports of the $q_{2j}$ are the disjoint triples of primary indices $4j, 4j + 1, 4j + 2$, with values $2, 1, 2$.
    Discard finitely many initial terms so that $v_k(x) = r$ on every remaining support and none of these supports lies over a prime dividing $n$.
    Every finite product of distinct remaining terms divides $x$ in $D$.
    For $n > 1$, all these terms are units modulo $nD$, by \eqref{eq:two-residues}.
    In their sequence of prefix products, some residue occurs infinitely often in the finite group $(D/nD)^{\times}$.
    Taking the nonempty blocks between successive occurrences gives disjoint blocks whose products $C_1, C_2, \ldots$ are each congruent to one modulo $nD$.
    For $n = 1$, take the remaining terms themselves as the $C_i$.
    Each $C_i$ belongs to $T$ and has a positive primary value, so it is a nonunit.
    Put $x_0 = x$ and $x_i = x/(C_1 \cdots C_i)$ for $i \ge 1$.
    We have $x_i \in D$, and $x = x_i C_1 \cdots C_i$ gives $x_i \equiv x \pmod{nD}$.
    Thus $x_i \in T$ and $x_{i - 1}T \subsetneq x_iT$ at every step.
    Factoring the $i$ nonunits $C_1, \ldots, C_i$ and the remaining factor $x_i$ in the atomic ring $T$ gives atomic factorizations of $x$ of length at least $i$.
\end{proof}

To finish Theorem~\ref{thm:subrings}, each $T_n$ contains $nt^2$, of order two, so it fails both ACCP and bounded factorization.
If $T$ has infinite additive index, it cannot contain any $mD$ with $m > 0$ because $D/mD$ is finite.
The conductor theorem therefore bounds the lengths of all decompositions of each of its nonzero elements into nonunits.
This gives bounded factorization and ACCP in $T$.
Every subring of $T$ also has infinite index in $D$, so $T$ is hereditarily a bounded factorization domain.
In particular this applies to every proper subring of $D$ containing $t$, since $t$ generates each additive quotient $D/T_n$.

For example, the element $t^2$ in $D$ has the factorization $t \cdot t$ and also factorizations of arbitrarily large length.
Exactly two factors in each have order one at $E$ by Proposition~\ref{prop:all-atomic}, and every additional factor has order zero.

\section{The exponent and the choice of paired primes}
\label{sec:sharpness}

The exponent in Theorem~\ref{thm:main} has an ideal-theoretic explanation.
In a Pr\"ufer domain whose Picard group is killed by two, squaring turns every factorization of an invertible ideal into a factorization of a principal ideal.
For an atom, this severely restricts which invertible ideals can contain it.

\begin{proposition}
    \label{prop:exponent-two}
    Let $R$ be a Pr\"ufer domain with $2\operatorname{Pic}(R) = 0$.
    Each atom generates either a radical ideal or the square of an invertible maximal ideal.
    If $R$ is atomic, it is a Dedekind domain or a field.
\end{proposition}

\begin{proof}
    Suppose $a$ is an atom and $aR$ is not radical.
    Choose $b \in R \setminus aR$ with $b^2 \in aR$.
    For instance, take $y \in \sqrt{aR} \setminus aR$ and let $b = y^{k - 1}$, where $k \ge 2$ is the least exponent with $y^k \in aR$.
    The ideal $I = (a, b)$ is invertible, properly contains $aR$, and satisfies $I^2 \subseteq aR$.
    Set $K = aI^{-1}$.
    Since $aR \subseteq I$, the ideal $K$ is integral, and $K = R$ would give $I = aR$.
    Multiplying $I^2 \subseteq aR$ by $aI^{-2}$ gives $aR \subseteq K^2$.
    The ideal $K^2$ is principal by the assumption on the Picard group, and it is proper since $K$ is proper.
    Irreducibility of $a$ forces $aR = K^2$.

    If $K$ were properly contained in a maximal ideal $N$, choose $z \in N \setminus K$ and put $J = K + zR$.
    The two invertible ideals $J$ and $KJ^{-1}$ are integral and proper, the latter because $K \subsetneq J$.
    Squaring $K = J(KJ^{-1})$ would express $aR$ as a product of two proper principal ideals.
    Thus $K$ is maximal, proving the first assertion.

    Now let $R$ be atomic and suppose some nonzero prime ideal $P$ is not finitely generated.
    It contains an atom, by factoring any nonzero element of $P$.
    Each atom in $P$ generates a radical ideal by the first assertion.
    Indeed, if $N^2 \subseteq P$ for an invertible maximal ideal $N$, primality gives $N \subseteq P$, so $P = N$ would be finitely generated.
    Fix an atom $a \in P$ and choose $b \in P \setminus aR$.
    The invertible ideal $J = (a, b)$ has principal square contained in $P$.
    Factor a generator of $J^2$ and choose an atom $y \in P$ occurring in this factorization.
    Then $J^2 \subseteq yR$ and $yR$ is radical, so $aR \subsetneq J \subseteq yR$.
    This makes $y$ a proper nonunit divisor of $a$, a contradiction.
    All prime ideals of $R$ are therefore finitely generated.
    Cohen's criterion \cite[Theorem~2]{Cohen} implies that $R$ is Noetherian.
    A Noetherian Pr\"ufer domain which is not a field is Dedekind.
\end{proof}

We close with a variation that distinguishes the examples by their finite-index subrings.
It retains exponent three and the same rational function field.

\begin{corollary}
    \label{cor:variation}
    There are continuum many pairwise nonisomorphic hereditarily atomic almost Dedekind SP-domains with fraction field $\mathbb{Q}(t)$, Picard group $\bigoplus_{j \ge 0} \mathbb{Z}/3\mathbb{Z}$, and an infinite strict principal-ideal chain starting at the square of an atom.
\end{corollary}

\begin{proof}
    Let $S$ be any infinite set of rational primes and assign the primary indices, in order, to its increasing enumeration.
    Keep a guard over every prime and omit the primary argument at $p \notin S$.
    The valuation at the first primary again gives $t$ value two.
    At a prime with just one argument, use $L_p(t) = t - \beta_p$ in Proposition~\ref{prop:construction} and use a constant $A_p$ to interpolate the single value of $t^{d_i} + A_p(t)$.
    The fixed-constant correction still preserves the prescribed value, and the slope in the derivative estimate remains nonzero since $L_p'(\beta_p) = 1$.
    For the coefficient estimate, $p^{2s}(t - \beta_p) = p^s(p^s t) - p^{2s}\beta_p$ still belongs to $A_s$, so the same normalization proves \eqref{eq:transfer-module}.
    The guard residue, Eisenstein, and geometric requirements are the same.
    These choices prove the construction proposition with one or two arguments as appropriate.

    The rational prime $p$ itself now isolates its sole guard when $p \notin S$.
    The proof of Lemma~\ref{lem:stalks} applies there with isolating value one.
    All primary coordinates and their relations persist, so Proposition~\ref{prop:ideals}, the chain, and Lemma~\ref{lem:extraction} give the same conclusions.
    The generic fibre proof uses the same geometric tests.
    At a one-argument prime the conductor generator divides $(t - \beta_p)^R$ and is divisible by $t - \beta_p$.
    The cancellation of a repeated factor in Lemma~\ref{lem:local-conductor} applies verbatim, giving conductor $(t - \beta_p)B'$.
    The constants in $T'$ already map onto $\mathbb{Z}_p$ under this one evaluation.
    Descent and the induction using the $a_n$ now prove hereditary atomicity as before.

    Write $D_S$ for this domain.
    At $p \notin S$, evaluation gives $D_S/pD_S \cong \mathbb{F}_p$ since its kernel is $pD_S$ and the constants map onto the residue field.
    At $p \in S$ the quotient is $\mathbb{F}_p^2$ by \eqref{eq:two-residues}.
    A subring of additive index $p$ must contain $pD_S$.
    The one-factor quotient has no proper unital subring, and the two-factor quotient has its diagonal of index $p$.
    Hence $S$ is exactly the set of primes at which $D_S$ has an index-$p$ subring.
    This is an invariant of the abstract ring, and there are continuum many infinite sets of primes.
\end{proof}


\begin{samepage}
    \paragraph{Authorship.}
    GPT-6.2 Astra discovered the result, building on earlier results discovered by GPT-6 Astra, Opus 5.5, and Fable 5.2.
    The proofs have received independent AI review.
    Human verification of the full construction is pending.
    \par
\end{samepage}

\begin{thebibliography}{9}
    \bibitem{Cohen}
    I.~S. Cohen, Commutative rings with restricted minimum condition,
    \emph{Duke Math. J.} \textbf{17} (1950), 27--42.
    \href{https://doi.org/10.1215/S0012-7094-50-01704-2}{doi:10.1215/S0012-7094-50-01704-2}.

    \bibitem{CGH}
    J.~Coykendall, F.~Gotti and R.~E.~Hasenauer, Hereditary atomicity in integral domains,
    \emph{J. Pure Appl. Algebra} \textbf{227} (2023), no.~4, 107249.
    \href{https://doi.org/10.1016/j.jpaa.2022.107249}{doi:10.1016/j.jpaa.2022.107249}.

    \bibitem{Grams}
    A.~Grams, Atomic rings and the ascending chain condition for principal ideals,
    \emph{Proc. Cambridge Philos. Soc.} \textbf{75} (1974), 321--329.
    \href{https://doi.org/10.1017/S0305004100048532}{doi:10.1017/S0305004100048532}.
\end{thebibliography}

\end{document}
